% !TeX program = xelatex
% !TeX encoding = UTF-8
\documentclass[11pt]{article}

\usepackage[a4paper,margin=1in]{geometry}
\usepackage{fontspec}
\usepackage{xeCJK}
\usepackage{amsmath}
\usepackage{amssymb}
\usepackage{mathtools}
\usepackage{booktabs}
\usepackage{tabularx}
\usepackage{enumitem}
\usepackage[hidelinks]{hyperref}

\IfFontExistsTF{Source Han Serif SC}{
  \setCJKmainfont{Source Han Serif SC}
}{
  \setCJKmainfont{Noto Serif CJK SC}
}
\IfFontExistsTF{Source Han Sans SC}{
  \setCJKsansfont{Source Han Sans SC}
}{
  \setCJKsansfont{Noto Sans CJK SC}
}
\setmainfont{TeX Gyre Termes}
\setsansfont{TeX Gyre Heros}
\xeCJKsetup{
  CJKspace=false,
  CJKecglue={},
  PunctStyle=kaiming
}

\linespread{1.18}
\setlength{\parindent}{2em}
\setlength{\parskip}{0.35em}
\setlength{\abovedisplayskip}{0.9em}
\setlength{\belowdisplayskip}{0.9em}
\setlength{\abovedisplayshortskip}{0.65em}
\setlength{\belowdisplayshortskip}{0.65em}
\setlength{\jot}{0.65em}
\renewcommand{\arraystretch}{1.18}
\allowdisplaybreaks
\numberwithin{equation}{section}

\DeclareMathOperator{\ConnectedComponents}{ConnectedComponents}

\title{Chemd 混合反应聚类数学设计}
\author{}
\date{}

\begin{document}

\maketitle

\section{设计目标}

本文描述 Chemd 反应聚类的数学设计。该设计的核心目标不是用单一文本相似度把反应粗略分组，而是先用化学证据约束可合并边界，再用 Chemd 文档图谱中的语义上下文补充排序、解释和候选关系。

设计上，任意两条反应 \(r_i\) 与 \(r_j\) 的关系由四类证据共同刻画：

\begin{enumerate}[label=\arabic*.]
  \item 反应指纹相似度 \(R_{ij}\)：描述反应整体结构变化是否相近。
  \item 反应向量相似度 \(X_{ij}\)：描述模型嵌入空间中的反应相似性。
  \item 反应中心相似度 \(C_{ij}\)：描述关键成键、断键或反应位点是否一致。
  \item 语义上下文相似度 \(S_{ij}\)：描述 Chemd 文档中的家族、步骤、条件、物料、路线、结果和观察是否相近。
\end{enumerate}

其中，\(R_{ij}\)、\(X_{ij}\) 与 \(C_{ij}\) 属于计算化学证据；\(S_{ij}\) 属于文档语义证据。严格反应簇必须由计算化学证据主导，语义上下文只用于增强排序和解释，不能单独覆盖冲突的化学证据。

\section{反应对相似度}

对任意两条反应 \(r_i\) 和 \(r_j\)，定义四个基础相似度分量：

\begin{equation}
\begin{aligned}
R_{ij} & = \text{反应指纹相似度}, \\
X_{ij} & = \text{反应向量相似度}, \\
C_{ij} & = \text{反应中心相似度}, \\
S_{ij} & = \text{语义上下文相似度}.
\end{aligned}
\end{equation}

\subsection{反应指纹相似度}

RDKit 反应指纹使用复合稀疏位向量。每一侧先拼接 path fingerprint 与 Morgan radius-2 fingerprint：

\begin{equation}
\begin{aligned}
A_i &= [A_i^{\mathrm{path}}, A_i^{\mathrm{morgan}}] \in \{0,1\}^{2048},\\
P_i &= [P_i^{\mathrm{path}}, P_i^{\mathrm{morgan}}] \in \{0,1\}^{2048},\\
G_i &= P_i \setminus A_i \in \{0,1\}^{2048},\\
L_i &= A_i \setminus P_i \in \{0,1\}^{2048},\\
F_i &= [A_i, P_i, G_i, L_i] \in \{0,1\}^{8192}.
\end{aligned}
\end{equation}

其中 \(A_i\) 是 reactant-side 指纹，\(P_i\) 是 product-side 指纹，\(G_i\) 是 product-gained 方向块，\(L_i\) 是 reactant-lost 方向块。任意两个位集合 \(U,V\) 的 Tanimoto 系数定义为：

\begin{equation}
T(U,V)
=
\frac{|U \cap V|}{|U \cup V|}.
\end{equation}

反应指纹相似度不是对 \(F_i\) 与 \(F_j\) 直接做一次整体 Tanimoto，而是按块加权：

\begin{equation}
R_{ij}
=
0.20\,T(A_i,A_j)
+0.20\,T(P_i,P_j)
+0.30\,T(G_i,G_j)
+0.30\,T(L_i,L_j).
\end{equation}

该定义保留 reactant side、product side、product-gained 和 reactant-lost 四类信息。新增与消失方向块权重更高，因此能区分氧化/还原这类正反方向相似但语义相反的反应；reactant/product 相似度提供背景支持，但不能单独覆盖反应中心冲突。

\subsection{反应向量相似度}

设 \(\mathbf e_i,\mathbf e_j \in \mathbb R^d\) 是两条反应的稠密向量表示。反应向量相似度使用截断后的余弦相似度：

\begin{equation}
X_{ij}
=
\max\left(
0,\,
\frac{\mathbf e_i \cdot \mathbf e_j}
{\|\mathbf e_i\|_2 \|\mathbf e_j\|_2}
\right).
\end{equation}

截断操作将负余弦值归零，避免把方向相反的向量解释为负证据。设计上，\(X_{ij}\) 适合作为结构指纹之外的补充信号，尤其用于捕获反应表示中的连续变化和模型学习到的潜在相似性。但它仍然不应单独决定严格聚类，因为向量相近可能来自上下文、底物类别或训练分布，而不一定代表反应中心一致。

\subsection{反应中心相似度}

反应中心相似度是离散化的化学判断信号：

\begin{equation}
C_{ij}
=
\begin{cases}
1.00, & \text{反应中心相同}, \\
0.72, & \text{反应中心兼容}, \\
0.00, & \text{反应中心冲突}, \\
\varnothing, & \text{反应中心不可用}.
\end{cases}
\end{equation}

这里的 \(1.00\) 表示关键成键、断键或反应位点在设计上可视为一致；\(0.72\) 表示两个反应中心并不完全相同，但仍可被视为同一类化学变化的兼容证据；\(0.00\) 表示存在明确冲突，后续硬门禁会优先处理这一情况；\(\varnothing\) 表示反应中心尚不可判定，不能被误当作低相似度。

\subsection{语义上下文相似度}

Chemd 文档图谱中的语义字段通常是集合型特征。对任意集合 \(A\) 和 \(B\)，使用 Jaccard 相似度：

\begin{equation}
J(A,B)
=
\frac{|A \cap B|}{|A \cup B|}.
\end{equation}

语义上下文由多个 Chemd 图谱特征加权混合得到：

\begin{equation}
\begin{aligned}
S_{ij}
=
& 0.25\,S^{\mathrm{family}}_{ij}
+ 0.20\,S^{\mathrm{procedure}}_{ij} \\
&+ 0.20\,S^{\mathrm{condition}}_{ij}
+ 0.15\,S^{\mathrm{material}}_{ij} \\
&+ 0.10\,S^{\mathrm{route}}_{ij}
+ 0.05\,S^{\mathrm{result}}_{ij} \\
&+ 0.05\,S^{\mathrm{observation}}_{ij}.
\end{aligned}
\end{equation}

各分量含义如下：

\begin{center}
\begin{tabularx}{\textwidth}{@{}llX@{}}
\toprule
分量 & 权重 & 设计含义 \\
\midrule
\(S^{\mathrm{family}}_{ij}\) & \(0.25\) & 反应家族或反应类型是否相近。 \\
\(S^{\mathrm{procedure}}_{ij}\) & \(0.20\) & 实验步骤、操作顺序或制备模式是否相近。 \\
\(S^{\mathrm{condition}}_{ij}\) & \(0.20\) & 温度、时间、气氛、溶剂、压力等条件是否相近。 \\
\(S^{\mathrm{material}}_{ij}\) & \(0.15\) & 底物、试剂、催化剂、添加剂等物料是否相近。 \\
\(S^{\mathrm{route}}_{ij}\) & \(0.10\) & 合成路线位置、上游来源或下游用途是否相近。 \\
\(S^{\mathrm{result}}_{ij}\) & \(0.05\) & 产率、选择性、转化率等结果描述是否相近。 \\
\(S^{\mathrm{observation}}_{ij}\) & \(0.05\) & 颜色、沉淀、放热、气体产生等观察记录是否相近。 \\
\bottomrule
\end{tabularx}
\end{center}

语义上下文的权重设计保持较低，是为了防止相同实验写法、相似路线描述或相同物料标签把化学上不相似的反应错误合并。它更适合回答“为什么这两条反应值得放在一起看”，而不是回答“这两条反应是否属于同一个严格化学簇”。

\section{混合相似度}

混合相似度 \(H_{ij}\) 使用带 RDKit 回补的加权策略。不可用值 \(\varnothing\) 不按 \(0\) 处理；当 RDKit 指纹相似度可用时，缺失证据的设计权重全部回补给 RDKit。这样可以让本地、确定性的 RDKit 指纹成为模型或 mapper 缺失时的主锚点，而不是让缺失的 RXNFP 或 reaction-center provider 扭曲总分。只有当 RDKit 不可用时，才退回到可用证据的归一化加权平均。

定义可用性指示函数：

\begin{equation}
\delta_Q(i,j)
=
\begin{cases}
1, & Q_{ij}\neq\varnothing, \\
0, & Q_{ij}=\varnothing.
\end{cases}
\end{equation}

定义四类证据的设计权重：

\begin{equation}
\begin{aligned}
w_R &= 0.30, &
w_X &= 0.25, &
w_C &= 0.25, &
w_S &= 0.20.
\end{aligned}
\end{equation}

设可用证据集合为

\begin{equation}
\mathcal A_{ij}
=
\{Q\in\{R,X,C,S\}\mid Q_{ij}\neq\varnothing\}.
\end{equation}

若 \(R_{ij}\neq\varnothing\)，有效权重定义为：

\begin{equation}
\begin{aligned}
\alpha_R
&=
w_R
+ w_X(1-\delta_X)
+ w_C(1-\delta_C)
+ w_S(1-\delta_S), \\
\alpha_X &= w_X\delta_X,\quad
\alpha_C = w_C\delta_C,\quad
\alpha_S = w_S\delta_S.
\end{aligned}
\end{equation}

此时混合相似度写为：

\begin{equation}
H_{ij}
=
\alpha_R R_{ij}
+ \alpha_X X_{ij}
+ \alpha_C C_{ij}
+ \alpha_S S_{ij}.
\end{equation}

若 \(R_{ij}=\varnothing\)，则不能凭空制造 RDKit 证据，退回到可用项归一化：

\begin{equation}
\begin{aligned}
H_{ij}
&=
\frac{N_{ij}}{D_{ij}}, \\
N_{ij}
&=
w_R R_{ij}\delta_R(i,j)
+ w_X X_{ij}\delta_X(i,j) \\
&\quad
+ w_C C_{ij}\delta_C(i,j)
+ w_S S_{ij}\delta_S(i,j), \\
D_{ij}
&=
w_R \delta_R(i,j)
+ w_X \delta_X(i,j)
+ w_C \delta_C(i,j)
+ w_S \delta_S(i,j).
\end{aligned}
\end{equation}

当 \(D_{ij}=0\) 时，表示该反应对没有任何可用证据。设计上令 \(H_{ij}=\varnothing\)，并且不生成相似边。该处理比默认给出 \(0\) 更保守，因为“没有证据”与“证据显示不相似”在聚类语义上并不相同。

计算化学证据是否存在定义为：

\begin{equation}
\begin{aligned}
\mathrm{Computed}_{ij}
&=
\mathbb 1[R_{ij}\neq\varnothing]
\lor
\mathbb 1[X_{ij}\neq\varnothing]
\lor
\mathbb 1[C_{ij}\neq\varnothing].
\end{aligned}
\end{equation}

\section{硬门禁}

严格聚类必须先通过硬门禁。硬门禁的作用是把“可解释的相似”分成两层：第一层判断化学上是否允许合并，第二层才使用混合分数判断是否足够接近。

反应中心冲突定义为：

\begin{equation}
\mathrm{CenterConflict}_{ij}
=
\mathbb 1[C_{ij}=0].
\end{equation}

当反应中心冲突且反应指纹也不足以支持相似时，直接拒绝：

\begin{equation}
\mathrm{HardReject}_{ij}
=
(C_{ij}=0)
\land
(R_{ij}<0.45).
\end{equation}

严格聚类允许条件定义为：

\begin{equation}
\begin{aligned}
\mathrm{StrictAllowed}_{ij}
=
&\ \mathrm{Computed}_{ij} \\
&\land \neg\mathrm{CenterConflict}_{ij} \\
&\land
\left(
R_{ij}\ge 0.70
\lor
X_{ij}\ge 0.85
\lor
C_{ij}=1.00
\right) \\
&\land H_{ij}\ge 0.75.
\end{aligned}
\end{equation}

该条件表达了三层设计意图：首先，严格簇必须有至少一种计算化学证据；其次，明确的反应中心冲突不能被语义相似度覆盖；最后，至少有一个强化学信号达到阈值，并且整体混合相似度足够高。换言之，语义上下文可以提高一个边的解释质量，但不能把化学证据不足的关系提升为严格聚类关系。

\section{边类型分类}

每个反应对最多生成一种边类型：

\begin{equation}
\begin{aligned}
E_{ij}
=
\begin{cases}
\mathrm{strict\_computed\_similarity},
& \begin{aligned}[t]
  &\mathrm{StrictAllowed}_{ij} \\
  &\land \neg\mathrm{HardReject}_{ij},
  \end{aligned} \\

\mathrm{candidate\_neighbor},
& \begin{aligned}[t]
  &\mathrm{Computed}_{ij} \\
  &\land H_{ij}\ge 0.55 \\
  &\land \neg\mathrm{StrictAllowed}_{ij},
  \end{aligned} \\

\mathrm{semantic\_group\_similarity},
& \begin{aligned}[t]
  &\neg\mathrm{Computed}_{ij} \\
  &\land S_{ij}\ge 0.65,
  \end{aligned} \\

\varnothing,
& \text{其他情况}.
\end{cases}
\end{aligned}
\end{equation}

三类边的设计语义如下：

\begin{enumerate}[label=\arabic*.]
  \item \(\mathrm{strict\_computed\_similarity}\)：进入严格反应簇的边，代表化学证据和混合分数同时满足合并条件。
  \item \(\mathrm{candidate\_neighbor}\)：有一定化学证据，但不足以直接合并的候选近邻，适合用于人工复核、推荐、局部扩展或后续更昂贵的相似度计算。
  \item \(\mathrm{semantic\_group\_similarity}\)：缺少计算化学证据时，仅由语义上下文形成的语义分组边，适合用于文档导航和主题组织，不应默认进入严格化学簇。
\end{enumerate}

\section{严格反应簇构造}

设 \(V\) 是所有反应节点集合。严格边集合定义为：

\begin{equation}
E_{\mathrm{strict}}
=
\{(i,j)\mid E_{ij}=\mathrm{strict\_computed\_similarity}\}.
\end{equation}

严格相似图为：

\begin{equation}
G_{\mathrm{strict}}
=
(V,E_{\mathrm{strict}}).
\end{equation}

严格反应簇定义为严格相似图的连通分量：

\begin{equation}
\mathcal C_{\mathrm{strict}}
=
\ConnectedComponents(G_{\mathrm{strict}}).
\end{equation}

候选近邻边和语义分组边保留为辅助结构，不参与严格反应簇的连通分量计算：

\begin{equation}
E_{\mathrm{candidate}}
=
\{(i,j)\mid E_{ij}=\mathrm{candidate\_neighbor}\},
\end{equation}

\begin{equation}
E_{\mathrm{semantic}}
=
\{(i,j)\mid E_{ij}=\mathrm{semantic\_group\_similarity}\}.
\end{equation}

这种分层构造可以避免“语义上相近但化学上不确定”的反应通过图连通性被间接并入严格簇。严格簇用于化学相似关系，候选近邻用于待确认关系，语义分组用于文档组织关系。

\section{簇置信度}

对任意严格反应簇 \(C\)，设 \(E_C\) 为簇内严格边集合。簇的平均相似度为：

\begin{equation}
\mathrm{cluster\_score}(C)
=
\frac{1}{|E_C|}
\sum_{(i,j)\in E_C} H_{ij}.
\end{equation}

簇内最弱边为：

\begin{equation}
\mathrm{min\_edge\_score}(C)
=
\min_{(i,j)\in E_C} H_{ij}.
\end{equation}

反应中心一致性为：

\begin{equation}
\begin{aligned}
\mathrm{center\_consistency}(C)
=
\frac{
|\{(i,j)\in E_C \mid C_{ij}\in\{1.00,0.72\}\}|
}{
|E_C|
}.
\end{aligned}
\end{equation}

簇置信度定义为：

\begin{equation}
\begin{aligned}
\mathrm{confidence}(C)
=
\begin{cases}
\mathrm{high},
& \begin{aligned}[t]
  &\mathrm{cluster\_score}(C)\ge 0.85 \\
  &\land\ \mathrm{min\_edge\_score}(C)\ge 0.75 \\
  &\land\ \mathrm{center\_consistency}(C)\ge 0.80,
  \end{aligned} \\

\mathrm{medium},
& \mathrm{cluster\_score}(C)\ge 0.70, \\

\mathrm{low},
& \text{其他情况}.
\end{cases}
\end{aligned}
\end{equation}

该置信度不是新的聚类条件，而是对已有严格簇质量的解释层。高置信度簇要求平均分、最弱边和反应中心一致性同时足够强；中置信度簇表示整体相似度可接受但仍可能存在弱边；低置信度簇需要在界面或后续分析中谨慎展示。

\section{设计解释}

该混合聚类设计可以概括为三条原则。

\begin{enumerate}[label=\arabic*.]
  \item \textbf{化学证据决定严格合并。}只有存在计算化学证据，并且没有明确反应中心冲突时，反应对才可能进入严格反应簇。
  \item \textbf{语义证据用于补充，而不是覆盖。}语义上下文可以提升排序质量、补充候选关系、解释簇画像，但不能单独把缺少化学支持的反应放入严格簇。
  \item \textbf{候选关系和语义分组必须与严格簇分离。}候选近邻适合用于人工复核、推荐和后续计算；语义分组适合用于文档导航和主题组织。二者都不应默认等价于严格化学聚类。
\end{enumerate}

因此，最终输出应至少区分三类结果：严格反应簇、候选近邻关系和语义分组。这种分层可以在保持化学可靠性的同时，利用 Chemd 文档结构提供更完整的解释和探索能力。

\section{建议编译方式}

本文档面向 Overleaf 的 XeLaTeX 编译链。请在 Overleaf 菜单中将编译器设置为 XeLaTeX。若项目字体列表中没有 \texttt{Source Han Serif SC}，本文档会退回到 \texttt{Noto Serif CJK SC}；该字体与思源宋体同源，适合作为 Overleaf 环境中的兼容替代。

\end{document}
